\chapter{ডেটা অভিধান}\label{app:dictionary}
\begin{table}[htbp]\centering\footnotesize
\caption{লেনদেন ডেটাসেটের ডেটা অভিধান}\label{tab:dictionary}
\begin{tabularx}{\linewidth}{@{}llX@{}}
\toprule
\textbf{কলাম} & \textbf{ধরন} & \textbf{বিবরণ} \\ \midrule
Transaction\_ID & টেক্সট & অনন্য লেনদেন শনাক্তকারী (যেমন DC2020000000) \\
Date & তারিখ & লেনদেনের তারিখ (dd-mm-yy) \\
Year, Month, Month\_Name, Quarter & পূর্ণসংখ্যা/টেক্সট & তারিখ থেকে নির্ণীত ক্যালেন্ডার-বৈশিষ্ট্য \\
Product\_ID, Product\_Name & টেক্সট & পণ্যের কোড ও নাম (\NProducts টি পণ্য) \\
Category & টেক্সট & পণ্যের ক্যাটাগরি (13টি) \\
Pack\_Size & টেক্সট & প্যাক সাইজ, যেমন 30g, 27--30g, 30g x 12, 5 pcs \\
Product\_Launch\_Year & পূর্ণসংখ্যা & পণ্য চালুর বছর \\
Observed\_Price\_BDT & পূর্ণসংখ্যা & মুদ্রিত সর্বোচ্চ খুচরা মূল্য \\
List\_Price\_BDT & পূর্ণসংখ্যা & বছরের কোম্পানি-নির্ধারিত তালিকা মূল্য \\
Market\_Price\_BDT & দশমিক & ডিসকাউন্টের আগে লেনদেন-মূল্য \\
Discount\_\% & দশমিক & ডিসকাউন্টের শতাংশ \\
Unit\_Selling\_Price\_BDT & দশমিক & ডিসকাউন্টের পরে একক মূল্য \\
Quantity & পূর্ণসংখ্যা & বিক্রীত ইউনিট \\
Gross\_Sales\_BDT & দশমিক & পরিমাণ $\times$ বাজার মূল্য \\
Discount\_BDT & দশমিক & ডিসকাউন্টের পরিমাণ \\
Net\_Sales\_BDT & দশমিক & মোট বিক্রয় $-$ ডিসকাউন্ট \\
Region & টেক্সট & বাংলাদেশের 11টি অঞ্চলের একটি \\
Channel & টেক্সট & General Trade, Modern Trade, Distributor, E-commerce \\
Customer\_Type & টেক্সট & Retailer, Super Shop, Distributor, Corporate, Online Customer \\ \bottomrule
\end{tabularx}
\end{table}

\chapter{টিউন করা Hyper-parameter}\label{app:hyper}
\begin{table}[htbp]
\centering
\footnotesize
\caption{Expanding-window যাচাইসহ randomised search-এ নির্বাচিত hyper-parameter (চাহিদা পূর্বাভাস)}
\label{tab:hyper}
\begin{adjustbox}{max width=\linewidth}
\begin{tabular}{l>{\raggedright\arraybackslash}p{10cm}}
\toprule
\textbf{Model} & \textbf{Selected hyper-parameters} \\
\midrule
Linear Regression & none (closed form) \\
Decision Tree & max\_depth=10, min\_samples\_leaf=40 \\
Random Forest & max\_depth=None, max\_features=0.33, min\_samples\_leaf=2, n\_estimators=500 \\
Gradient Boosting & learning\_rate=0.05, max\_depth=2, n\_estimators=200, subsample=1.0 \\
XGBoost & colsample\_bytree=0.6, learning\_rate=0.02, max\_depth=3, min\_child\_weight=10, n\_estimators=600, reg\_lambda=1, subsample=0.9 \\
K-Nearest Neighbors & n\_neighbors=20, weights=distance \\
Support Vector Regression & C=0.3, epsilon=0.05, gamma=0.01 \\
\bottomrule
\end{tabular}
\end{adjustbox}
\end{table}

\chapter{উদাহরণ SQL কোয়েরি}\label{app:sql}
নিচের কোয়েরিগুলো স্টার স্কিমায় (\cref{fig:star}) চলে এবং \cref{ch:eda}-এর প্রধান সংখ্যাগুলো সরাসরি SQL-এ পুনরুৎপাদন করে। এগুলো একটি
Power BI প্রতিবেদনের ডেটা-উৎস হিসেবেও ব্যবহার করা যায়।
\lstinputlisting[language=SQL]{sql/kpi_queries.sql}

\chapter{প্রকল্পটি পুনরুৎপাদন ও প্রকাশ}\label{app:reproduce}
\section{প্রকল্পের কাঠামো}
\begin{lstlisting}
dan_cake_project/
  run_pipeline.py          one command for the whole framework
  src/dancake/             config, data, eda, features, forecasting, classification,
                           clustering, anomaly, insights, export, website
  dashboard/app.py         Streamlit dashboard
  data/  processed/ warehouse/ dashboard/   generated data
  outputs/ metrics/ tables/ models/          generated results
  report/  main.tex, main_bn.tex, chapters/, chapters_bn/, figures/, tables/, tables_bn/
  docs/                    project website (GitHub Pages)
  sql/  tests/  guides/
\end{lstlisting}

\section{নিজের কম্পিউটারে চালানো}
\begin{lstlisting}[language=bash]
python -m venv .venv
.venv\Scripts\activate            # Windows  (source .venv/bin/activate on macOS/Linux)
pip install -r requirements.txt
python run_pipeline.py            # about 20 minutes
python -m pytest tests            # automated tests
streamlit run dashboard/app.py    # open http://localhost:8501
\end{lstlisting}
বাংলা প্রতিবেদন XeLaTeX দিয়ে কম্পাইল করতে হয়: \texttt{tectonic -X compile main\_bn.tex}, অথবা Overleaf-এ Compiler = XeLaTeX।

\section{বিনামূল্যে প্রকাশ}
\begin{enumerate}
  \item একটি বিনামূল্যের GitHub অ্যাকাউন্ট ও একটি পাবলিক রিপোজিটরি তৈরি করে প্রকল্পটি push করুন (\texttt{.gitignore} কাঁচা CSV বাদ দেয়;
        ছোট ড্যাশবোর্ড-ডেটা অন্তর্ভুক্ত থাকে)।
  \item \textbf{ড্যাশবোর্ড:} GitHub দিয়ে Streamlit Community Cloud-এ সাইন ইন করে \emph{New app} বেছে নিন, রিপোজিটরি, branch
        \texttt{main} ও ফাইল \texttt{dashboard/app.py} নির্বাচন করে deploy করুন। অ্যাপটি \texttt{.streamlit.app}-এ শেষ হওয়া একটি ঠিকানা পাবে।
  \item \textbf{ওয়েবসাইট:} রিপোজিটরির settings-এ \emph{Pages} খুলে \emph{Deploy from a branch}, branch \texttt{main}, folder
        \texttt{/docs} বেছে নিন। সাইটটি \path{https://<user>.github.io/<repo>}-এ প্রকাশিত হবে।
  \item প্রকৃত লিংকসহ ওয়েবসাইটটি আবার তৈরি করুন (\texttt{src} ফোল্ডারের ভেতর থেকে), তারপর commit ও push করুন:
\begin{lstlisting}[language=bash]
python -m dancake.website --dashboard-url <app url> \
                          --repo-url <repo url>
\end{lstlisting}
\end{enumerate}

\section{পরবর্তী প্রকল্পে কাঠামোটির পুনঃব্যবহার}
নতুন একটি বিক্রয় ডেটাসেটের জন্য কেবল \texttt{src/dancake/config.py} (পথ, কলামের নাম, train/test তারিখ) এবং ফিচার-ফাংশনের কলাম-তালিকা
বদলাতে হবে। কাঠামোটি (যাচাই, পরিষ্কারকরণ, ওয়্যারহাউস, EDA, leakage-মুক্ত ফিচার, বেঞ্চমার্কসহ মডেল, ড্যাশবোর্ড, অন্তর্দৃষ্টি ও স্বয়ংক্রিয়
প্রতিবেদন) অপরিবর্তিতভাবে মজুদ, অর্থ বা গ্রাহক অ্যানালিটিক্সের মতো অন্য ক্ষেত্রেও প্রয়োগ করা যায়।
